\documentclass[12pt]{article}
\usepackage[T1]{fontenc}
\usepackage[utf8]{inputenc}
\usepackage[brazil]{babel}
\usepackage[margin=1in]{geometry}
\usepackage{ragged2e}
\usepackage[normalem]{ulem}
\setlength{\parindent}{0pt}
\setlength{\parskip}{0.8\baselineskip}
\begin{document}
\justifying
\noindent Verifica-se que a questão objeto da controvérsia jurídica suscitada no recurso manejado pela parte recorrente esteve afetada no representativo da controvérsia - \textbf{Tema n. }\textbf{183} da TNU - (PU no processo n. PEDILEF 0500796-67.2017.4.05.8307/PE),  afetado em 21/06/2018 (afetado pelo Pleno da TNU) foi julgado (trânsito em julgado em 24/09/2019 (no STF - RE 1194635/PE)), por meio do qual foi elucidada a seguinte questão:

\noindent \textit{"Decidir se INSS tem responsabilidade civil pelos danos patrimoniais ou morais decorrentes de empréstimo consignado não autorizado."}

\noindent Restou firmada a seguinte tese:

\noindent \textit{"I - O INSS não tem responsabilidade civil pelos danos patrimoniais ou extrapatrimoniais decorrentes de “empréstimo consignado”, concedido mediante fraude, se a instituição financeira credora é a mesma responsável pelo pagamento do benefício previdenciário, nos termos do art. 6º, da Lei n. 10.820/03; II – O INSS pode ser civilmente responsabilizado por danos patrimoniais ou extrapatrimoniais, se demonstrada negligência, por omissão injustificada no desempenho do dever de fiscalização, se os “empréstimos consignados” forem concedidos, de forma fraudulenta, por instituições financeiras distintas daquelas responsáveis pelo pagamento dos benefícios previdenciários. A responsabilidade do INSS, nessa hipótese, é subsidiária em relação à responsabilidade civil da instituição financeira." (Tema 183 TNU).}

\noindent Com efeito, depreende-se da leitura do voto condutor do acórdão que o mesmo se encontra em conformidade com o decidido no  tema pela TNU.

\noindent A proposito, consta do acordao hostilizado: (...) Empréstimo consignado, Bancários, Contratos de Consumo, DIREITO DO CONSUMIDOR

\noindent Nessas condições, o conteúdo do Acórdão recorrido é consentâneo com o entendimento da TNU, firmado em julgamento de recurso representativo de controvérsia, o que atrai a incidência da regra prevista pelo art. 14, III, \textit{b,} da Resolução CJF n. 586/2019:

\noindent \textit{\textbf{Art. 14.}}\textit{ Decorrido o prazo para contrarrazões, os autos serão conclusos ao magistrado responsável pelo exame preliminar de admissibilidade, que deverá, de forma sucessiva: (...) }\textit{\textbf{III }}\textit{- negar seguimento a pedido de uniformização de interpretação de lei federal interposto contra acórdão que esteja em conformidade com entendimento consolidado: }\textit{\textbf{b)}}\textit{ em }\uline{\textit{\textbf{recurso representativo de controvérsia pela Turma Nacional de Uniformização}}}\textit{ ou em pedido de uniformização de interpretação de lei dirigido ao Superior Tribunal de Justiça;(...).}

\noindent Posto isso, com arrimo no \uline{\textbf{art. 14, III, }}\uline{\textit{\textbf{b}}}\textit{,} da Resolução CJF n. 586/2019 c/c art. 1.030, I, do CPC, \textbf{NEGO SEGUIMENTO }\textbf{a}o\textbf{ PU}\textbf{.}

\noindent Intimem-se. Aguarde-se o decurso do prazo recursal. Após, certifique-se o trânsito em julgado e devolva-se ao juízo de origem.
\end{document}
